\chapter{Introduction}
\label{ch:intro}

A supernumerary robotic limb is an extra arm worn on the body, working
alongside the wearer's own limbs rather than replacing or rehabilitating
them~\cite{yang2021srlreview,prattichizzo2021augmentation}. What separates
such a device from a fixed manipulator is consequential for control: the base
of the robot is a person. It accelerates, sways and steps, and it cannot be
told to hold still.

This thesis reports work on Imperial College London's Dr.\ Octopus platform
(the MUlti-limb Virtual Environment~\cite{muve}), which mounts
two Kinova Gen3 seven-degree-of-freedom arms~\cite{kinova2022gen3} on a
backpack frame and commands them from a \emph{second} person: an operator who
cannot feel the arms, controlling them for a wearer who cannot command them.
No prior work has measured both what shared autonomy buys such an operator
and what it costs the person carrying the arms, in one system, under the
same trial (\cref{sec:gap}); that gap is what this project is built to
close.

\begin{figure}[htbp]
  \centering
  \includegraphics[width=0.8\linewidth]{photos/platform_scene_hd.png}
  \caption[The platform as built]{The platform as built: two Kinova Gen3
  arms on a backpack frame, on a torso mannequin, photographed by the scene
  camera. The two teal columns are the laboratory's fixed manipulators and
  are not part of this system.}
  \label{fig:platform}
\end{figure}

Two premises shaped the engineering: position is taken from whichever
sensor observes each component best rather than from the slave's full
kinematic chain, and a body-worn master \emph{will} lose channels over its
life, so channel loss is treated as an operating condition that is
announced, never papered over -- on a manipulator beside a person's head, a
fabricated position is worse than none.

\section{Background}
\label{ch:background}

\textbf{Supernumerary robotic limbs.} Parietti and Asada's founding
demonstrations brace a limb against the environment to stabilise the
wearer's posture~\cite{parietti2014bracing,parietti2013dynamic}, which turns
the wearer's body from a disturbance source into part of a closed chain.
Later reviews map the design space along consistent axes -- mounting
location trades reach against the moment arm imposed on the wearer,
actuation trades bandwidth against worn mass -- and note that the field
reports embodiment far more often than task
performance~\cite{yang2021srlreview,prattichizzo2021augmentation,tong2021srlreview,zhang2023srlreview,zheng2026quantitative}.
Worn mass binds hardest: each Gen3 masses \SI{8.2}{\kilo\gram}, so the pair
with harness exceeds \SI{17}{\kilo\gram}, and simulated wearer loads suggest
the biomechanical cost is not negligible~\cite{moon2024impact}.

\textbf{The operator and the wearer as different people.} Fusion mounts a
robotic head and two arms on a worn frame, driven by a remote
operator~\cite{saraiji2018fusion} -- but its three modes are levels of
\emph{physical coupling}, not machine autonomy, and the publication reports
no user study. Work on near-body interaction found that raising
\emph{apparent} autonomy (a concealed operator, a wearer with no task of
their own) raised physiological arousal and lowered self-reported
trust~\cite{zhou2026proxemics}: the wearer's experience is not a free
variable, though this is not evidence about real autonomy or an occupied
wearer. The two-person case also differs mechanically: one-person
base-motion compensation reaches \SI{1.37 \pm 0.58}{\milli\metre} of
end-effector error~\cite{zhang2024motioncomp}, but there the disturbance
comes from the operator's own nervous system, which supplies an efference
copy of it; here it comes from somebody else's, which supplies none.

\textbf{Master interfaces for dissimilar kinematics.} Guggenheim and Asada
argue for miniature physical masters on the grounds that geometric
similarity supplies feedback inherently~\cite{guggenheim2021haptic}; the
argument covers similarity of \emph{form}, not \emph{scale}, and this
system's master-to-slave reach ratio is roughly one to three. Joint-space
correspondence preserves posture but needs every master joint measured
accurately; task-space correspondence decouples the chains at the cost of
the slave's redundant degree of freedom.

\textbf{Shared autonomy.} Dragan and Srinivasa formalise arbitration as
policy blending, weighted by prediction confidence~\cite{dragan2013policyblending};
Javdani et al.\ recast it as a POMDP over the user's goal and show that
assistance which does not wait for high confidence helps operators finish
faster~\cite{javdani2015hindsight}. Neither was developed for a robot mounted
on a second person, where the cost of assisting toward the wrong goal is not
a slower trial but a manipulator moving toward someone who did not command
it.

\textbf{Measuring embodiment, workload and trust.} Ownership over an
artificial limb is measurable and decomposes into ownership, agency and
location, which can move independently~\cite{botvinick1998rubber,longo2008embodiment}.
Workload separates into the NASA Task Load Index's six
subscales~\cite{hart1988nasatlx}; physical exertion is better tracked with a
Borg scale~\cite{borg1982perceived}, appropriate for an uncompensated master
arm or a carried mass. Feedback that improves control does not reliably
improve ownership~\cite{pinardi2023feedback}, so performance and embodiment
must be measured separately -- which is why the pilot here measures the
former and claims nothing about the latter.

\section{The gap}
\label{sec:gap}

No study measures the operator and the wearer together, under the same
autonomy, in the same trial: shared autonomy has been shown to help an
operator where the robot is attached to nobody, and raising apparent
autonomy has been shown to lower a wearer's trust where the autonomy was a
concealed human. Until both are measured in one system, the exchange rate
between what shared autonomy buys the operator and what it costs the person
carrying the arms is unknown. The full review, with every citation's scope
stated precisely, is in \cref{app:background}.

\section{Aim, objectives and hypotheses}
\label{sec:aim}

The aim was to build and verify, on the platform of \cref{fig:platform}, a
two-person teleoperation system on which that exchange rate can be
measured, and to begin measuring it. Six objectives:

\begin{enumerate}
  \item Design and build an instrumented mannequin master reproducing the
    Gen3 morphology at reduced scale, characterised against recorded data.
  \item Derive a master-to-slave command mapping that degrades predictably as
    sensing channels are lost, and quantify what each loss costs.
  \item Implement a safety architecture appropriate to a manipulator mounted
    beside a person, validated by fault injection rather than inspection.
  \item Characterise the reachable workspace of the mounted configuration and
    determine which manipulation tasks it can support.
  \item Establish a verification methodology producing traceable evidence for
    every claim, and apply it throughout.
  \item Run a pilot comparison of direct control against shared-autonomy
    assistance with real operators, and identify what stands between it and
    the full two-person study the architecture was designed for.
\end{enumerate}

The study the system is built for states five hypotheses, with direction,
before any data (\cref{app:study}): task completion time is shorter under
shared autonomy (H1); operator workload is lower (H2); wearer-reported trust
is \emph{lower} under shared autonomy (H3), against the direction of H1 and
H2; the wearer's sense of agency falls while ownership does not (H4); and
cross-arm interference falls while phase lag does not (H5). The pilot
reported here tests H1 directly; no workload instrument was administered, so
H2 is untested, and H3--H5 need the separate wearer it did not have.

Each objective is reported in \cref{ch:resultsmain} with its evidence class
-- implemented, verified in simulation, measured on hardware, or measured on
a pilot session. The planning report's haptic-master objective was not
delivered (\cref{ch:discussionmain,app:development}). \textbf{No result
below the pilot was validated on physical hardware carrying a load.}
